\chapter{Methods}
\label{ch:methods}

Having covered the necessary prerequisites in \Cref{ch:vpt}, we now present our
own method. We describe our modifications to VPT, the architecture and training
of the model, and how these components come together in the final solution.

\section{Direct and indirect illumination}

\begin{figure}[t]
    \centering
    \begin{tikzpicture}
    \filldraw[draw=cyan!15, fill=cyan!5, thick, xscale=1.9]
        (4.406,0.408)
        .. controls (4.095,0.173)  and (3.695,0.032)  .. (3.297,0.0)
        .. controls (2.899,-0.032) and (2.504,0.048)  .. (2.063,0.068)
        .. controls (1.619,0.089)  and (1.131,0.05)   .. (0.751,0.202)
        .. controls (0.372,0.356)  and (0.104,0.702)  .. (0.0,1.058)
        .. controls (-0.104,1.417) and (-0.046,1.786) .. (0.053,2.128)
        .. controls (0.153,2.471)  and (0.296,2.784)  .. (0.564,2.977)
        .. controls (0.835,3.169)  and (1.229,3.239)  .. (1.6,3.29)
        .. controls (1.967,3.344)  and (2.31,3.377)   .. (2.735,3.425)
        .. controls (3.163,3.471)  and (3.672,3.53)   .. (4.035,3.37)
        .. controls (4.395,3.211)  and (4.606,2.833)  .. (4.774,2.454)
        .. controls (4.938,2.073)  and (5.056,1.69)   .. (5.0,1.332)
        .. controls (4.947,0.975)  and (4.721,0.645)  .. (4.406,0.408)
        -- cycle;

    \coordinate (sun) at (5,5);
    \coordinate (start) at (-0.8,1.9);

    \node[circle, fill, inner sep=1.6pt] (a) at (1.4,1.9) {};
    \node[circle, fill, inner sep=1.6pt] (b) at (4,0.8) {};
    \node[circle, fill, inner sep=1.6pt] (c) at (8,1.4) {};
    \node[circle, fill, inner sep=1.6pt] (d) at (6,2.7) {};
    
    \shade[inner color=yellow!60, outer color=white] (sun) circle (0.8cm);
    \node[circle, draw=yellow!90, fill=yellow!20, thick,
        minimum size=0.6cm, inner sep=0pt] (sun_node) at (sun) {};
    \draw[-Stealth, semithick] (start) -- (a);

    % Direction radiance
    \draw[-Stealth, dashed, semithick] (a) -- (sun_node);

    % Indirect radiance 
    \draw[-Stealth, dotted, semithick] (a) -- (b);
    \draw[-Stealth, dotted, semithick] (b) -- (c);
    \draw[-Stealth, dotted, semithick] (c) -- (d);
    \draw[-Stealth, dotted, semithick] (d) -- (sun_node);
\end{tikzpicture}
    \caption{
        Paths for generating a sample (solid) and computing direct (dashed)
        and indirect (dotted) radiance for this sample.
    }
    \label{fig:scattering}
\end{figure}

\begin{figure}[b]
    \centering
    \begin{subfigure}{0.3\textwidth}
        \centering
        \includegraphics[width=0.85\textwidth]{global.png}
    \end{subfigure}
    \begin{subfigure}{0.3\textwidth}
        \centering
        \includegraphics[width=0.85\textwidth]{direct.png}
    \end{subfigure}
    \begin{subfigure}{0.3\textwidth}
        \centering
        \includegraphics[width=0.85\textwidth]{indirect.png}
    \end{subfigure}
    
    \caption{
        Result of composing global illumination (left) from direct (middle)
        and indirect (right) radiances.
    }
    \label{fig:illumination}
\end{figure}

In most related work, only the indirect radiance component is cached or
approximated by a neural network, while direct radiance is computed exactly using
Monte Carlo integration. We use the same approach. Because VPT does not natively 
support separating these components, we rework its path tracing logic. Following
\cite{bauer_photon_2024}, we split the computation into three separate shaders:
\begin{itemize}
    \item \emph{Sample generation.} Perform delta tracking until a scattering
    or an absorption event occurs. This gives us a sample in which we hit
    the volume.
    \item \emph{Direct radiance.} Starting at the generated sample, perform
    delta tracking until we either hit a light or the photon is absorbed. In
    the case of scattering, the ray is discarded because it is not regarded as
    single scattering anymore.
    \item \emph{Indirect radiance.} As with direct radiance, perform delta
    tracking but allow scattering events and discard rays with fewer than two
    bounces, i.e., real scattering events, when hitting a light.
\end{itemize}
\Cref{fig:scattering} illustrates these shaders for a single sample. Once the
direct and indirect radiance have been computed, a separate shader \emph{composes}
them into global illumination by taking their average, weighted by the number of 
direct and indirect samples. When the \emph{accumulate} option is enabled, all
three illumination components (global, direct, and indirect) are additionally
accumulated over successive frames, each independently of the others. During evaluation, this design made it possible to produce images of all three
components, visible in \Cref{fig:illumination}, in a single experiment and to 
switch between them without restarting the computation.

We also rework the logic for handling iterations. Previously, if the maximum
number of bounces was not reached within the current iteration, VPT would
continue the path in the next iteration, i.e., the next frame. We change this
so that each path is computed from start to finish within a single iteration.
This change allows us to introduce an additional parameter, \emph{samples},
which controls the number of photons traced per pixel.

\section{Ground truth}

As we already compute indirect radiance when rendering the image, a separate
pass to generate ground truth data for the model is not needed. Reusing these
values means the model will learn to predict the radiance values needed for the
current camera angle. Since there is enough data available in each frame, we
store only the first valid sample out of \emph{samples} rays sent per pixel.

For each sample, we store the value either when an absorption event occurs, in
which case the radiance is zero, or when we hit the scene boundary after at
least two real scattering events. In the latter case, the radiance is the 
transmittance multiplied by the environment sample, which is uniform white light
in our setup, so the final stored value is just the transmittance.

For each pixel, we store eight values in a \ac{gpu} buffer: the position
$(x, y, z)$, the sample’s direction azimuth and elevation $(\varphi, \theta)$,
and the radiance $(r, g, b)$. Note that the position and direction are taken
from the start of the path, i.e., from the sample generated by the sample
generation shader, not from the last ray position. For a $512 \times 512$ image,
this amounts to $\qty{8}{\mega\byte}$ of data per frame.

Before this data can be used, it must be copied from the \ac{gpu} to main memory.
We do this using asynchronous mapping in WebGPU. This copying can take multiple
frames, in which case we drop data until the copying is finished. After the data
is available, we either add it to a ZIP archive as a binary file or send it to our
training server. Using a ZIP archive, we can download $\qty{4}{\giga\byte}$ of
data, which is useful during development and later for repeatable evaluation
experiments.

\section{Radiance Field Network}

Our neural network model, \emph{Radiance Field Network}, is based on the
Photon Field Network model from~\cite{bauer_photon_2024}. We modify the model
for our input data and for running it in a browser. Because our data are
generated for a specific renderer configuration, we drop the coefficient $g$,
which is constant. Position and direction are still encoded using
Multiresolution Hash Encoding from~\cite{muller_instant_2022}, but we remove the
log transform of the output because we do not observe any improvement in our
results. Instead, we clamp the output to the interval $[0, 1]$. The
architecture diagram is shown in \Cref{fig:arch}.

The position and direction vectors have separate encoding blocks. Each block
contains $L$ levels, and each level outputs an $F$-dimensional embedding. The
levels span resolutions from $N_{\min}$ to $N_{\max}$ as follows:
\begin{align}
    b &= \exp\left( \frac{\log N_{\max} - \log N_{\min}}{L - 1} \right), \\ 
    N_l &= \left\lfloor N_{\min} \cdot b^l \right\rfloor.
\end{align}
On each level, we first find the corners of the voxel containing the input
vector. These corners are hashed to obtain indices, which are then used to
retrieve embeddings from the table with maximum size $T$. These embeddings are
linearly interpolated in proportion to the distance between the input vector
and the voxel corners. All embeddings are then concatenated.

In our case, we use
$L = 11$,
$F = 4$,
$N_{\min}^{\text{pos}} = 8$, $N_{\max}^{\text{pos}} = 128$,
$N_{\min}^{\text{dir}} = 16$, $N_{\max}^{\text{dir}} = 64$,
and $T = 2^{18}$. All of these values, along with $3$ fully connected layers
of size $16$ to produce the final output, were selected using an automatic 
Bayesian hyperparameter search.

\begin{figure}[t]
    \centering
    \begin{tikzpicture}[
    enc/.style={
        rounded corners, draw=violet!30, fill=violet!10,
        align=center, font=\linespread{0.8}\selectfont,
        minimum width=5, inner sep=6, very thick},
    fc/.style={
        rounded corners, draw=cyan!30, fill=cyan!8,
        align=center, font=\linespread{0.8}\selectfont,
        minimum width=80, inner sep=6, very thick},
    relu/.style={
        rounded corners, draw=orange!40, fill=orange!8,
        minimum width=80, inner sep=4, very thick},
    box/.style={rounded corners, draw=gray!45, fill=gray!12,
        inner sep=4, very thick, font=\small},
    pad/.style={rounded corners, inner sep=6, very thick},
    arr/.style={-Stealth, thick},
    line/.style={-, thick}
]
    % MLP ---------------------------------------------------------------
    \node[fc] (fc1) at (0,0) {Fully\\connected};
    \node[relu, below=0.15cm of fc1] (relu1) {ReLU};
    
    \node[fc, below=0.7cm of relu1] (fc2) {Fully\\connected};
    \node[relu, below=0.15cm of fc2] (relu2) {ReLU};

    \node[fc, below=0.7cm of relu2] (fc3) {Fully\\connected};
    % \node[relu, below=0.15cm of fc3] (relu3) {ReLU};

    \node[pad, draw=olive!40, fill=olive!10, minimum width=5,
        below=0.7cm of fc3] (clamp) {Clamp};

    \node[below=1cm of clamp] (out) {$(r, g, b)$};

    \draw[line] (fc1) -- (relu1);
    \draw[arr] (relu1) -- (fc2);
    \draw[line] (fc2) -- (relu2);
    \draw[arr] (relu2) -- (fc3);
    \draw[line] (fc3) -- (clamp);
    % \draw[arr] (relu3) -- (clamp);
    \draw[arr](clamp) -- (out);

    % Encoding ----------------------------------------------------------
    \node (xyz)    at (-7.5,0 |- out) {$(x,y,z)$};
    \node (angles) at (-5,0 |- out) {$(\varphi,\theta)$};

    \coordinate (mid) at ($(xyz)!0.4!(angles)$);
    \node[enc, minimum width=6.5cm, minimum height=5cm, anchor=south]
        (box2) at ([xshift=0.15cm]mid |- clamp.south) {};
    \node[enc, minimum width=6.5cm, minimum height=5cm, anchor=south]
        (box1) at ([shift={(-0.3cm,0.3cm)}]box2.south) {};
    \node[draw=violet!30, fill=violet!15, rounded corners,
        minimum width=5.15cm, minimum height=4.4cm,
        anchor=east, very thick] (level)
        at ([xshift=-0.3cm]box1.east) {};
    \node[anchor=east] at ([xshift=-0.5mm]level.west) {$L\times$};

    % Right column
    \node[box, anchor=south east, minimum width=3.2cm] (voxels)
        at ([shift={(-0.2cm,0.2cm)}]level.south east) {Voxel corners};
    \node[box, anchor=south east, minimum width=2.3cm] (hash)
    at ([yshift=0.4cm]voxels.north east) {Hashing};
    \node[box, anchor=south east, minimum width=2.3cm] (emb)
        at ([yshift=0.4cm]hash.north east) {Embedding}; 

    % Left column
    \node[box, minimum width=2.2cm] (weights) at (xyz |- emb) {Weights};

    % Top
    \node[box, anchor=south east, minimum width=4.7cm] (int)
        at ([yshift=0.4cm]emb.north east) {Linear interpolation}; 

    % Arrows
    \draw[arr] (xyz) |- (voxels);
    \draw[arr] (xyz.north) -- (xyz.north |- weights.south);
    \draw[Stealth-, thick] (hash.south) -- (hash.south |- voxels.north);
    \draw[arr] (hash.north) -- (hash.north |- emb.south);
    \draw[arr] (weights.north) -- (weights.north |- int.south);
    \draw[arr] (emb.north) -- (emb.north |- int.south);
    \draw[arr] ([xshift=-1.3cm]voxels.north) --
        ([xshift=-1.3cm]voxels.north |- weights.south);

    \node[pad, draw=purple!50, fill=purple!20,
        minimum width=4cm, anchor=south] (cat) at ([yshift=7.4cm]mid) {Concat};

    \draw[line] (angles.north) |- (angles.north |- box1.south); 
    \draw[arr] ([xshift=-0.7cm]cat.south |- int.north) 
        -- ([xshift=-0.7cm]cat.south);
    \draw[arr] ([xshift=0.7cm]cat.south  |- box1.north)
        -- ([xshift=0.7cm]cat.south);
    \draw[arr] (cat) |- (fc1);
\end{tikzpicture}
    \caption{Radiance Field Network architecture.}
    \label{fig:arch}
\end{figure}

The model is optimized with the Adam optimizer, using a learning rate of $0.01$
and an MSE loss function. Because the model is intended for use in an online
setting, we train it for a single epoch.

\section{Training server}

To use our model during rendering, we have to train it in an online setting.
Offline training would require waiting until the model is fully trained, whereas
online training makes it possible to start generating images while the model is
still being trained in the background. There are not many deep learning libraries
available in the browser, and most of them do not support training. Those that do 
usually only have experimental support for WebGPU. We could implement 
backpropagation from scratch, but because browser performance falls short of 
desktop performance, we instead decided to implement a WebSocket server that runs
training in PyTorch.

\begin{figure}[!b]
    \centering
    \begin{tikzpicture}[
    rgt/.style={-{Stealth}, semithick, shorten <=4pt, shorten >=4pt},
    lft/.style={{Stealth}-, semithick, shorten <=4pt, shorten >=4pt},
    lbl/.style={midway, above, font=\small},
]
    \node[rounded corners, draw=gray!40, fill=gray!10,
        minimum width=2cm, inner sep=6, thick] (b) at (0,0) {Browser};
    \node[rounded corners, draw=gray!40, fill=gray!10,
        minimum width=2cm, inner sep=6, thick] (s) at (7,0) {Server};
        
    \draw[dashed, gray!60, semithick] (b.south) -- ++(0,-7.3);
    \draw[dashed, gray!60, semithick] (s.south) -- ++(0,-7.3);
    
    \draw[rgt] (b.south |- 0,-1) -- (s.south |- 0,-1)
        node[lbl] {WebSocket connect};
    \draw[lft] (b.south |- 0,-2) -- (s.south |- 0,-2)
        node[lbl] {\textsf{model-created}};

    \draw[rgt] (b.south |- 0,-3.75) -- (s.south |- 0,-3.75)
        node[lbl] {\textsf{ground-truth(pixel data)}};
    \draw[lft] (b.south |- 0,-4.75) -- (s.south |- 0,-4.75)
        node[lbl] {\textsf{metrics(val loss, train time)}};

    \draw[gray!60, semithick]
        ([xshift=-5mm] b.south |- 0,-2.75) rectangle
        ([xshift=5mm] s.south |- 0,-5.25);
    \node[draw=gray!60, fill=gray!10, semithick, font=\small\bfseries,
          inner sep=3pt, outer sep=0pt, anchor=north west]
        at ([xshift=-5mm] b.south |- 0,-2.75) {loop};

    \draw[rgt] (b.south |- 0,-6.3) -- (s.south |- 0,-6.3)
        node[lbl] {\textsf{model-request}};
    \draw[lft] (b.south |- 0,-7.3) -- (s.south |- 0,-7.3)
        node[lbl] {\textsf{model-weights(bytes)}};
\end{tikzpicture}
    \caption{
        Sequence diagram of message passing between the browser and the
        training server during online training. When the connection is
        established, a new model is created. Received ground truth data are
        used for training, while the final model weights are obtained using a
        separate message type.
    }
    \label{fig:server}
\end{figure}

The main protocol between the browser and the server is illustrated in
\Cref{fig:server}. When the browser establishes a new connection to the server,
a new model instance is initialized. After the browser receives confirmation
that the model is ready, the user can enable the \emph{train} option to start 
sending ground truth data. This data is sent on every frame, except when data is
still being copied. Instead of storing the bytes in a ZIP archive as before, they 
are now sent directly to the server with the \emph{ground-truth} header. The
server first parses and preprocesses the data. It removes empty pixels, invalid 
angles, and out-of-bounds coordinates, while normalizing the angles to the unit
interval. A single training step is performed, and the validation loss is returned 
to be displayed in the browser. When the user enables the \emph{predict} option,
a \emph{model-request} message is sent to the server. The server responds by 
creating a ZIP archive of the model weights, which is sent as bytes in a message
to the browser.

Model training can be restarted either by reconnecting to the server or by using
a separate \emph{model-reset} message type. This allows us to restart training if
the model gets stuck in a poor local minimum. To prevent this during evaluation,
we added support for setting the seed using a special message sent at the start
of evaluation to make it reproducible. Model resetting is also performed each
time the VPT configuration is changed, since the model is trained for a specific
configuration.

\section{Bilateral filter}

Since our ground truth data values are generated from the camera view, we can
imagine them as forming a 2D image. This is only an approximation, because the
generated samples do not hit the volume at the same depth, meaning the actual
image plane is not flat. However, this still allows us to apply some image
filtering methods to our ground truth data before handing it off for training.
In this case, we propose that the model can learn this filtering, which could be
beneficial if the filtering algorithm is computationally more demanding.

In our approach, we use the \emph{bilateral filter} for this use
case~\cite{paris_gentle_2007}. It is similar to Gaussian filtering but preserves
edges by weighting neighbors not only by their distance but also by their color
difference. It has three parameters: $\sigma_k$, $\sigma$, and a threshold $\tau$.
For each pixel with color $\mathbf{c}$, we consider neighbors with colors
$\mathbf{c}_n$ in the neighborhood with radius
\begin{equation}
    r = \sigma_k \cdot \sigma.
\end{equation}
If the pixel is not a background pixel, we compute the spatial weight
\begin{equation}
    w_d = \exp\left(-\frac{dx^2 + dy^2}{2\sigma^2}\right),
\end{equation}
and the range weight
\begin{equation}
    w_r = \exp\left(-\frac{\norm{\mathbf{c} - \mathbf{c}_n}}{2\tau^2}\right).
\end{equation}
We then multiply these weights and use the result to compute the weighted average
over the entire neighborhood. A comparison between the filtered and unfiltered
images is shown in \Cref{fig:filter}.

\begin{figure}[t]
    \centering
    \begin{subfigure}{0.46\textwidth}
        \centering
        \includegraphics[width=0.85\textwidth]{filter_off.png}
    \end{subfigure}
    \begin{subfigure}{0.46\textwidth}
        \centering
        \includegraphics[width=0.85\textwidth]{filter_on.png}
    \end{subfigure}
    \caption{
        Comparison of an image rendered without filtering (left) and
        with the bilateral filter enabled (right).
    }
    \label{fig:filter}
\end{figure}

\section{Neural render}

All components described in the previous sections come together to form the
algorithm we call \emph{neural render}. It is based on path tracing, with
indirect radiance computation replaced by neural network inference.

Path tracing is described in \Cref{alg:pt}. A shader first generates sample
points, which are then used to compute both direct and indirect radiance. If
filtering is enabled, we apply it to both. If frame accumulation is disabled,
previous radiance values are discarded; otherwise, they are updated with the
newly computed values. The final image depends on the current display mode
(global illumination, direct radiance, or indirect radiance). After rendering
is complete, the resulting indirect radiance values are sent to the training
server.

The neural render described in \Cref{alg:nr} follows similar steps. We still
generate sample points and compute direct radiance in the same way, but now the
indirect radiance is computed using a neural network, so the volume data is no
longer needed. If filtering is enabled, the model is already trained on filtered
data, so we only have to filter the direct radiance. The remaining steps stay
the same, except for sending data to the training server. Because we no longer
compute indirect radiance values, we do not have ground truth values to send to
the server.

The model forward pass is implemented in a WebGPU shader, so the data stays on
the device for the entire iteration. Before neural rendering starts, the model
weights received from the server are copied to buffers on the \ac{gpu}. We
process all samples for a single pixel in the same shader invocation. To simulate
transparency, which is the result of delta tracking in path tracing, we introduce
an additional rejection of some samples. We toss a coin and reject samples with
probability $P_n(\mathbf{x})$, which is already computed and stored in the sample
generation shader.

\begin{algorithm}
    \caption{Path tracing}
    \label{alg:pt}
    \DontPrintSemicolon
    \KwIn{%
        volume $\mathcal{V}$,
        camera $\mathcal{C}$,
        configuration $\gamma$,
        radiance $L$}
    \KwOut{rendered image $I$, updated radiance $L$}
    \BlankLine
    
    $\mathbf{X}, \boldsymbol\Omega \leftarrow$ GenerateSamplePoints(%
        $\mathcal{V}$, $\mathcal{C}$, $\gamma$)\;
    $L_d \leftarrow$ ComputeDirectIllumination(%
        $\mathcal{V}$, $\gamma$, $\mathbf{X}$, $\boldsymbol\Omega$)\;
    $L_i \leftarrow$ ComputeIndirectIllumination(%
        $\mathcal{V}$, $\gamma$, $\mathbf{X}$, $\boldsymbol\Omega$)\;
    \BlankLine

    \If{filtering enabled}{
        $L_d, L_i \leftarrow$ ApplyBilateralFilter($L_d$, $L_i$, $\gamma$)\;
    }
    \If{not accumulate}{
    $L \leftarrow \mathbf{0}$\;
    }
    \BlankLine
    
    $L \leftarrow$ AccumulateRadiance($L$, $L_d$, $L_i$)\;
    $I \leftarrow$ ComposeImage($L$, $\gamma$) \;
    TrainServerSend(\texttt{"ground-truth"}, $L_i$)\;
    \BlankLine
    
    \Return{$I$, $L$}
\end{algorithm}

\begin{algorithm}
    \caption{Neural render}
    \label{alg:nr}
    \DontPrintSemicolon
    \KwIn{%
        volume $\mathcal{V}$,
        camera $\mathcal{C}$,
        configuration $\gamma$, radiance $L$,
        model $\mathcal{N}$}
    \KwOut{rendered image $I$, radiance $L$}
    \BlankLine

    $\mathbf{X}, \boldsymbol\Omega \leftarrow$ GenerateSamplePoints(%
        $\mathcal{V}$, $\mathcal{C}$, $\gamma$)\;
    $L_d \leftarrow$ ComputeDirectIllumination(%
        $\mathcal{V}$, $\gamma$, $\mathbf{X}$, $\boldsymbol\Omega$)\;
    $L_i \leftarrow$ PredictIndirectIllumination(%
        $\mathcal{N}$, $\mathbf{X}$, $\boldsymbol\Omega$)\;
    \BlankLine
    
    \If{filtering enabled}{
        $L_d \leftarrow$ ApplyBilateralFilter($L_d$, $\gamma$)\;
    }
    \If{not accumulate}{
    $L \leftarrow \mathbf{0}$\;
    }
    \BlankLine

    $L \leftarrow$ AccumulateRadiance($L$, $L_d$, $L_i$)\;
    $I \leftarrow$ ComposeImage($L$, $\gamma$)\;
    \BlankLine
    
    \Return{$I$, $L$}
\end{algorithm}
